
      \documentclass[fleqn]{article}      
      \usepackage{amsmath}
      \usepackage{fontspec}
      \usepackage{kotex} % 한국어 지원  

      \begin{document}      
      쉬운 문제: \\\\
1. 행렬의 정의는 무엇인가요? \\\\
2. \(2 \times 2\) 행렬의 예를 제시하세요. \\\\
3. 행렬의 덧셈이란 무엇인가요? \\\\
4. 다음 행렬을 더하세요: \(\begin{bmatrix} 1 & 2 \\ 3 & 4 \end{bmatrix} + \begin{bmatrix} 5 & 6 \\ 7 & 8 \end{bmatrix}\) \\\\
5. 행렬의 크기란 무엇을 의미하나요? \\\\

중간 문제: \\\\
1. 두 행렬의 곱셈에 필요한 조건은 무엇인가요? \\\\
2. 다음 행렬을 곱하세요: \(\begin{bmatrix} 1 & 2 \\ 3 & 4 \end{bmatrix} \times \begin{bmatrix} 5 & 6 \\ 7 & 8 \end{bmatrix}\) \\\\
3. 다음 행렬의 전치 행렬을 구하세요: \(\begin{bmatrix} 1 & 2 \\ 3 & 4 \end{bmatrix}\) \\\\
4. 행렬의 스칼라 배율이란 무엇인가요? \\\\
5. \( \begin{bmatrix} 1 & 0 \\ 0 & 1 \end{bmatrix} \) 행렬의 성질을 설명하세요. \\\\

어려운 문제: \\\\
1. 다음 행렬의 역행렬을 구하세요: \(\begin{bmatrix} 1 & 2 \\ 3 & 4 \end{bmatrix}\) \\\\
2. 두 행렬이 같다는 조건을 설명하세요. \\\\
3. 다음 행렬의 행렬식(Determinant)을 구하세요: \(\begin{bmatrix} 1 & 2 \\ 3 & 4 \end{bmatrix}\) \\\\
4. 행렬의 고유값과 고유벡터란 무엇인가요? \\\\
5. 행렬의 선형 독립성에 대해 설명하세요. \\\\

객관식 문제: \\\\
쉬움: \\\\
1. 행렬의 크기를 나타내는 것은? \\\\
   a) 행과 열의 수 \\\\
   b) 숫자의 합 \\\\
   c) 숫자의 곱 \\\\
  
2. 다음 중 행렬의 덧셈이 가능한 경우는? \\\\
   a) \(2 \times 2\) 행렬과 \(3 \times 3\) 행렬 \\\\
   b) \(2 \times 2\) 행렬과 \(2 \times 2\) 행렬 \\\\
   c) \(1 \times 2\) 행렬과 \(2 \times 1\) 행렬 \\\\
  
3. 다음 중 단위 행렬의 정의는? \\\\
   a) 모든 원소가 1인 행렬 \\\\
   b) 대각선 원소가 1이고 나머지는 0인 행렬 \\\\
   c) 모든 원소가 0인 행렬 \\\\
  
중간: \\\\
1. 다음 중 행렬의 곱셈에 대한 설명으로 올바른 것은? \\\\
   a) 행과 열의 수가 같아야 곱할 수 있다. \\\\
   b) 첫 번째 행렬의 열의 수와 두 번째 행렬의 행의 수가 같아야 곱할 수 있다. \\\\
   c) 두 행렬의 크기가 같아야 곱할 수 있다. \\\\
  
2. 다음 중 전치 행렬의 성질로 올바른 것은? \\\\
   a) 전치 행렬은 대칭 행렬이다. \\\\
   b) \((A^T)^T = A\)이다. \\\\
   c) 전치 행렬의 행렬식은 0이다. \\\\
  
3. 행렬의 스칼라 배율에서 스칼라가 0일 때의 결과는? \\\\
   a) 단위 행렬 \\\\
   b) 제로 행렬 \\\\
   c) 대칭 행렬 \\\\
  
어려움: \\\\
1. 다음 중 행렬의 역행렬에 대한 설명으로 올바른 것은? \\\\
   a) 역행렬은 항상 존재한다. \\\\
   b) \(AA^{-1} = I\)이다. \\\\
   c) 역행렬은 항상 대칭 행렬이다. \\\\
  
2. 행렬식이 0인 경우를 설명하세요. \\\\
   a) 행렬이 선형 독립일 때 \\\\
   b) 행렬이 선형 종속일 때 \\\\
   c) 행렬이 대칭일 때 \\\\
  
3. 다음 중 고유값의 성질로 올바른 것은? \\\\
   a) 모든 행렬은 고유값을 가진다. \\\\
   b) 고유값은 항상 정수이다. \\\\
   c) 고유값은 행렬의 행렬식과 관련이 있다. \\\\
  
      
      \end{document} 
  